% Part VI -- Algorithms, Optimization, and Machine Learning
% Chapter 16: Algorithmic thinking and efficiency

\h{Algorithmic Thinking and Efficiency}
\label{ch:16}

Programming syntax tells Python what to execute. An \emph{algorithm} is the
finite, ordered method that explains how a problem will be solved. Two programs
can produce the same answer while requiring very different amounts of work.
This chapter develops the habits used throughout computer science: specify the
problem, trace state, argue correctness, count important operations, and choose
a data structure that supports the algorithm.

\begin{NoteBox}
\textbf{Prerequisites.} You should understand lists, indexing, loops,
conditionals, functions, exceptions, and the idea of a loop invariant from
\booklink{ch:05}{Chapter 5}.
\end{NoteBox}

\begin{tcolorbox}[title=Learning outcomes,colframe=orange,colback=white,arc=2mm]
After completing this chapter, you should be able to:
\begin{itemize}
    \item distinguish a problem specification, an algorithm, and a program;
    \item state preconditions, postconditions, and useful edge cases;
    \item trace linear and binary search by hand;
    \item explain why binary search requires sorted input;
    \item use a loop invariant as a concise correctness argument;
    \item compare constant, logarithmic, linear, and quadratic growth;
    \item explain an educational sorting algorithm and use Python's stable
          \c{sorted()} function in applications; and
    \item validate two implementations against the same contract.
\end{itemize}
\end{tcolorbox}

\hh{Separate the problem, algorithm, and program}

Suppose the problem is ``find the position of a requested room identifier.'' A
complete \emph{specification} defines the input, expected output, and behavior
when the identifier is absent. An \emph{algorithm} describes the ordered search
strategy. A Python \emph{program} expresses that strategy using variables,
loops, comparisons, and returns.

\begin{SyntaxBox}{Write the contract before the loop}
For the search functions in this chapter:
\begin{itemize}
    \item the input is a sequence \c{values} and one \c{target};
    \item the output is the first matching index, or \c{-1} if absent;
    \item inputs may be empty; and
    \item binary search additionally requires ascending sorted order.
\end{itemize}
The contract is language-independent. It lets us compare algorithms without
changing the question they answer.
\end{SyntaxBox}

\begin{ExplainBox}{Correctness, efficiency, and implementation are separate claims}
Correctness asks whether every permitted input produces an output that satisfies
the contract. Efficiency asks how required time or storage grows as the input
grows. Implementation quality asks whether the Python expression is readable,
testable, and appropriate for its environment. A fast program can be wrong; a
correct algorithm can be implemented inefficiently; and two correct programs
can differ greatly in clarity. Evaluate each claim with its own evidence.
\end{ExplainBox}

\hh{Use test cases that represent different paths}

One successful example is not evidence that an algorithm is correct. A small
search test set should include the first item, a middle item, the last item, a
missing target, repeated values, and an empty sequence.

\begin{lstlisting}[
    language=python,
    style=block,
    caption={Describe search cases before implementing the algorithm},
    label={c16_search_cases},
    captionpos=b
]
test_cases = [
    ([4, 9, 13], 4, 0),
    ([4, 9, 13], 9, 1),
    ([4, 9, 13], 13, 2),
    ([4, 9, 13], 8, -1),
    ([], 8, -1),
]
\end{lstlisting}

Each tuple stores one input sequence, one target, and the expected index. The
cases become reusable evidence when several implementations share the same
contract.

\hh{Linear search: inspect candidates in order}

Linear search begins at index zero and stops at the first match. It works on
sorted or unsorted sequences.

\begin{lstlisting}[
    language=python,
    style=block,
    caption={Return the first matching index with linear search},
    label={c16_linear_search},
    captionpos=b
]
def linear_search(values, target):
    for index, value in enumerate(values):
        if value == target:
            return index
    return -1

print(linear_search([18, 24, 15, 31], 15))
print(linear_search([18, 24, 15, 31], 20))
\end{lstlisting}

\begin{SyntaxBox}{An early return ends a successful search}
\c{enumerate(values)} supplies both the current index and value. The equality
comparison asks one question per candidate. When a match occurs, \c{return}
immediately ends the function. The final \c{return -1} is reached only after
the loop has ruled out every position.
\end{SyntaxBox}

\begin{ExplainBox}{Linear-search invariant}
At the beginning of iteration \c{index}, the target has not appeared in any
earlier position. The comparison either finds it at the current position or
extends that fact by one position. If the loop ends, every position has been
ruled out, so returning \c{-1} satisfies the contract.
\end{ExplainBox}

\hh{Binary search: discard half of a sorted range}

Sorted order provides additional information. Compare the target with the
middle item. If the middle item is too small, every item on its left is also
too small. If it is too large, every item on its right is also too large.

\begin{lstlisting}[
    language=python,
    style=block,
    caption={Search an ascending sequence by repeatedly halving the range},
    label={c16_binary_search},
    captionpos=b
]
def binary_search(sorted_values, target):
    low = 0
    high = len(sorted_values) - 1

    while low <= high:
        mid = (low + high) // 2
        candidate = sorted_values[mid]

        if candidate == target:
            return mid
        if candidate < target:
            low = mid + 1
        else:
            high = mid - 1

    return -1
\end{lstlisting}

\bookfigure[0.98\textwidth]{Figures/Generated/ch16_binary_search_trace.pdf}
{Binary search for 36 compares a middle value, discards the impossible half, and preserves only the range that can still contain the target.}
{fig:ch16-binary-search}

\begin{SyntaxBox}{Track the inclusive search interval}
\begin{itemize}
    \item \c{low} and \c{high} are the first and last possible indices.
    \item \c{// 2} produces an integer midpoint index.
    \item \c{low = mid + 1} removes the midpoint and everything below it.
    \item \c{high = mid - 1} removes the midpoint and everything above it.
    \item \c{low > high} means no possible index remains.
\end{itemize}
The invariant is: if the target exists, it lies somewhere from \c{low} through
\c{high}, inclusive. Every update must preserve that statement.
\end{SyntaxBox}

\begin{TryBox}{Expose the precondition}
Run binary search on \c{[18, 4, 31, 9]} without sorting. A returned answer may
be wrong even though the code raises no exception. Sort a copy first with
\c{sorted(values)}, run again, and explain why an algorithm can be syntactically
valid while violating its contract.
\end{TryBox}

\hh{Compare efficiency by order of growth}

Exact runtime depends on hardware, Python version, input layout, and many other
details. Algorithm analysis first asks how the number of important operations
grows as input size $n$ grows.

\bookfigure[0.98\textwidth]{Figures/Generated/ch16_algorithm_growth.pdf}
{Constant, logarithmic, linear, and quadratic operation counts respond very differently as input size grows; the table makes the orders comparable at larger values.}
{fig:ch16-growth}

\begin{center}
\begin{tabular}{lll}
\textbf{Growth} & \textbf{Notation} & \textbf{Typical example} \\
Constant & $O(1)$ & read one list item by index \\
Logarithmic & $O(\log n)$ & binary search \\
Linear & $O(n)$ & linear search \\
Quadratic & $O(n^2)$ & compare many pairs in nested loops \\
\end{tabular}
\end{center}

Big-O notation describes an upper-order growth pattern, not a stopwatch time.
Linear search may be perfectly appropriate for a short unsorted list. Binary
search becomes valuable when the data are already sorted or will be searched
many times. Sorting once also has a cost, so the complete workflow matters.

\hh{Sorting reveals nested-loop reasoning}

Selection sort is useful for learning because its invariant is visible: before
iteration \c{start}, every earlier position contains its final smallest value.
The algorithm scans the remaining suffix, finds its smallest element, and swaps
that element into position.

\begin{lstlisting}[
    language=python,
    style=block,
    caption={Implement selection sort for study},
    label={c16_selection_sort},
    captionpos=b
]
def selection_sort(values):
    result = values.copy()

    for start in range(len(result)):
        smallest = start

        for index in range(start + 1, len(result)):
            if result[index] < result[smallest]:
                smallest = index

        result[start], result[smallest] = (
            result[smallest],
            result[start],
        )

    return result
\end{lstlisting}

The outer loop chooses the next output position. The inner loop scans the
remaining candidates, producing quadratic growth. The parallel assignment
swaps two values without losing either one.

\begin{NoteBox}
\textbf{Use the standard tool in real applications.} Python's \c{sorted()}
returns a new list and uses a highly engineered stable sorting algorithm. The
educational implementation above is for tracing and complexity analysis, not
for replacing the standard library.
\end{NoteBox}

\hh{Sort records with a key function}

Algorithms often need a rule for comparing structured records. A key function
returns the comparison value for each item.

\begin{lstlisting}[
    language=python,
    style=block,
    caption={Sort room records by area without changing the originals},
    label={c16_sorted_key},
    captionpos=b
]
rooms = [
    {"name": "Studio", "area_m2": 31.0},
    {"name": "Office", "area_m2": 18.0},
    {"name": "Lobby", "area_m2": 24.0},
]

by_area = sorted(rooms, key=lambda room: room["area_m2"])
print(by_area)
\end{lstlisting}

\c{lambda room: room["area_m2"]} is a small anonymous function. For each room,
it returns the numeric field used as the sort key. Stability means records with
equal keys preserve their earlier relative order. The official Python
\href{https://docs.python.org/3/howto/sorting.html}{Sorting HOWTO} explains
keys, stability, and multi-level sorting.

\hh{Applied example: validate and compare two searches}
\label{sec:c16-application}

The complete companion program is
\c{examples/c16_search_algorithms.py}. It gives both algorithms the same
contract, counts comparisons, asserts that they return the same index, and
plots operation counts for growing inputs.

\begin{lstlisting}[
    language=python,
    style=block,
    caption={Count comparisons without changing the search result},
    label={c16_counted_search},
    captionpos=b
]
def linear_search(values, target):
    comparisons = 0
    for index, value in enumerate(values):
        comparisons += 1
        if value == target:
            return index, comparisons
    return -1, comparisons


def binary_search(sorted_values, target):
    low = 0
    high = len(sorted_values) - 1
    comparisons = 0

    while low <= high:
        mid = (low + high) // 2
        comparisons += 1
        candidate = sorted_values[mid]

        if candidate == target:
            return mid, comparisons
        if candidate < target:
            low = mid + 1
        else:
            high = mid - 1

    return -1, comparisons


room_ids = [104, 112, 119, 125, 131, 144, 152, 167, 181]
target_id = 167

linear_result = linear_search(room_ids, target_id)
binary_result = binary_search(room_ids, target_id)

print("Linear:", linear_result)
print("Binary:", binary_result)
assert linear_result[0] == binary_result[0]
\end{lstlisting}

\begin{SyntaxBox}{Instrumentation should preserve the contract}
The functions now return a tuple containing the answer and a measurement.
Incrementing \c{comparisons} immediately before each equality test defines what
is being counted. The \c{assert} is an executable claim that both algorithms
locate the same target. If measurement code changes the algorithm's answer, the
experiment is invalid.
\end{SyntaxBox}

\hh{Common mistakes and useful diagnoses}

\begin{itemize}
    \item \textbf{Binary search receives unsorted data.} Validate the
          precondition or sort a copy when that cost is appropriate.
    \item \textbf{The interval never shrinks.} Update past \c{mid} with
          \c{mid + 1} or \c{mid - 1}; otherwise the loop may repeat forever.
    \item \textbf{Only successful cases are tested.} Include absent targets and
          empty inputs.
    \item \textbf{Big-O is treated as an exact duration.} It describes growth,
          not milliseconds on one computer.
    \item \textbf{A teaching sort replaces \c{sorted()}.} Use the standard
          implementation unless the learning goal is the algorithm itself.
\end{itemize}

\hh{Practice ladder}

\hhh{Level 0 -- Trace}

Trace binary search for target 9 in \c{[4, 9, 13, 18, 22, 27, 31]}. Record
\c{low}, \c{high}, \c{mid}, and the candidate after every comparison.

\hhh{Level 1 -- Modify}

Change both search functions to return \c{None} rather than \c{-1} when the
target is absent. Update the tests so the contract remains explicit.

\hhh{Level 2 -- Explain}

Explain why binary search is $O(\log n)$ and selection sort is $O(n^2)$ without
referring to their Python syntax.

\hhh{Level 3 -- Build}

Write \c{is_sorted(values)} that returns \c{True} for an empty or one-item
sequence and otherwise checks every neighboring pair.

\hhh{Level 4 -- Challenge}

Modify binary search to return the first index of a repeated target. State the
new invariant and test cases before writing the loop.

\hh{Practice solutions}

\textbf{Level 0:} The midpoints contain 18, then 9; the second comparison finds
the target.

\textbf{Level 2:} Binary search halves the remaining candidates, so doubling
the input adds approximately one comparison. Selection sort repeatedly scans a
large remaining suffix, producing a total proportional to the square of the
input size.

\textbf{Level 3}

\begin{lstlisting}[
    language=python,
    style=block,
    caption={One adjacent-pair sortedness check},
    label={c16_solution_sorted},
    captionpos=b
]
def is_sorted(values):
    for index in range(1, len(values)):
        if values[index] < values[index - 1]:
            return False
    return True
\end{lstlisting}

The loop begins at index one because every comparison needs a previous item.
The function returns early at the first inversion and returns \c{True} only
after ruling out every inversion.

\hh{Chapter checkpoint}

\begin{enumerate}
    \item How does a specification differ from an algorithm and a program?
    \item What is the additional precondition for binary search?
    \item State the linear-search invariant.
    \item Why must binary search update past the midpoint?
    \item What do $O(1)$, $O(\log n)$, $O(n)$, and $O(n^2)$ compare?
    \item Why is selection sort useful educationally but not the normal Python choice?
    \item What does a sorting key function return?
\end{enumerate}

\hh{Checkpoint answers}

\begin{enumerate}
    \item A specification defines the required behavior, an algorithm defines
          the method, and a program expresses that method in a language.
    \item The sequence must be sorted according to the same comparison rule.
    \item Before each comparison, no earlier position contains the target.
    \item The compared midpoint is no longer a possible candidate.
    \item How important operation counts grow with input size.
    \item It exposes invariants and nested-loop growth; \c{sorted()} is more
          capable and efficient for applications.
    \item The comparison value extracted from one record.
\end{enumerate}

\begin{tcolorbox}[title=Chapter summary,colframe=orange,colback=white,arc=2mm]
Algorithmic thinking begins with a contract and edge cases, continues through
a state trace and correctness invariant, and ends with an efficiency argument.
Syntax implements the method, but the method is the transferable knowledge.
These habits prepare you to formulate search problems over continuous values
in \booklink{ch:17}{Chapter 17}.
\end{tcolorbox}
